\section{Research Plan}

\subsection{Scope of the November 2026 minimum viable benchmark}

The target for 30 November 2026 is a tested Python library rather than a
single experiment script. Its minimum viable release will (i) sample valid
discrete-time formulas, (ii) execute them to generate reproducible datasets,
(iii) derive several explanation ground truths from the retained formulas,
(iv) accept forecasting models through a small adapter interface, and
(v) produce comparable forecasting and explanation reports. The first release
will prioritize deterministic autoregressive, cross-variable, nonlinear, and
irrelevant-input mechanisms. Regime switching and stochastic mechanisms will
enter the release only after the deterministic path passes its validation
tests, because they introduce additional questions about state labels,
innovation terms, and the object to be explained.

The first stochastic generator will use additive zero-mean Gaussian
innovations with a sampled scale. This initial restriction makes the expected
error distribution explicit and testable. Additional innovation families,
including heavy-tailed, heteroscedastic, and regime-dependent noise, will be
considered after the Gaussian generator and its evaluation interface pass the
validation gates.

The intended dependency order is
\[
\boxed{
\begin{gathered}
\text{formula AST}
\longrightarrow
\text{executable generator}
\longrightarrow
\text{time-series dataset}
\\[1mm]
\Downarrow
\\[1mm]
\text{explanation ground truths}
\longrightarrow
\text{model adapters}
\longrightarrow
\text{evaluation}
\end{gathered}
}.
\]
Consequently, model experiments do not begin before both the observations and
their references can be generated and checked. A lightweight software
interface is specified early so that the generator does not become tied to one
forecasting package, but the interface is exercised on real models only after
the ground-truth layer is complete.

\subsection{Formula sampler}

Following the expression-first strategy of \citet{dascoli2022deep}, the
sampler will first create an abstract syntax tree (AST) and only then execute
it. Let $\mathcal E$ denote a valid symbolic expression, either the complete
right-hand side of a recurrence or one of its subexpressions. A minimal
recursive grammar is
\[
\begin{aligned}
\mathcal E ::= {}& c
&& \text{sampled constant,}\\
&\mid \operatorname{LAG}(j,\ell)
&& \text{series $j$ observed $\ell$ steps earlier,}\\
&\mid \operatorname{ADD}(\mathcal E_1,\mathcal E_2)
&& \text{sum of two valid expressions,}\\
&\mid \operatorname{MUL}(\mathcal E_1,\mathcal E_2)
&& \text{product of two valid expressions,}\\
&\mid \operatorname{UNARY}_{\varphi}(\mathcal E_1)
&& \text{unary function applied to one expression,}
\end{aligned}
\qquad \ell\geq 1.
\]
The symbol $\mathcal E$ is therefore a grammar non-terminal: it is a
placeholder for any expression that can be constructed by repeatedly applying
these rules. The symbols $\mathcal E_1$ and $\mathcal E_2$ are child
expressions constructed by the same grammar, $c$ is a sampled constant, $j$
identifies a source series, and $\varphi$ is selected from a controlled set of
safe unary functions.

For example, the symbolic tree
\[
\begin{aligned}
\mathcal E={}&\operatorname{ADD}\!\bigl(
    \operatorname{MUL}(0.7,\operatorname{LAG}(1,1)),\\
&\hspace{38mm}\operatorname{SIN}(\operatorname{LAG}(1,3))
\bigr)
\end{aligned}
\]
is executed as the recurrence
\[
x_{1,t}=0.7x_{1,t-1}+\sin(x_{1,t-3}).
\]
Extensions will add explicit interaction, regime, multiscale, and innovation
nodes. The AST is retained alongside every series as the canonical description
of the mechanism; a rendered equation is a human-readable view of that same
object.

The sampler must control tree depth, operator counts, admissible domains,
maximum lag, coefficient ranges, and duplicated or inactive branches. Candidate
formulas are rejected when short pilot simulations diverge, collapse to an
almost constant series, produce non-finite values, or fail to use the sampled
dependencies. Each accepted formula is serialized with its seed and complete
parameter record so that it can be reconstructed exactly.

\subsection{Series generation and dataset contract}

For target series $k$, the AST is compiled into an executable recurrence
\[
x_{k,t}=g_{k,1}\!\left(x_{:,t-1},\ldots,x_{:,t-L};\theta\right)
       +\varepsilon_{k,t}.
\]
Initial conditions are sampled first, a burn-in interval is discarded, and the
remaining observations are divided chronologically into training, validation,
and test windows. The generator records the formula, coefficients, initial
conditions, innovations, regime sequence when applicable, seeds, window
boundaries, and temporal convention. This metadata distinguishes the process
that generated the observations from the input window later presented to a
model.

Multi-step reference outputs are produced by recursively executing the known
generator. Thus $g_{k,h}$ denotes the output at horizon $h$ after unrolling the
one-step recurrence $h$ times. The benchmark will expose the same sample through
a canonical batch containing observed history, optional known future
covariates, target values, horizon, masks, and metadata. An adapter may transform
that batch into a package-specific dataframe or tensor, but it cannot silently
change which information the model receives.

\subsection{Explanation ground truths}

Only after an accepted formula and its observations exist does the library
derive explanation references. One generated example can expose four
compatible views:

\begin{enumerate}
    \item a binary source--lag support extracted from the AST and, for
    $h>1$, from the recursively unrolled dependency graph;
    \item signed coefficients for mechanisms whose effect is represented by a
    coefficient, together with realized coefficient--input terms;
    \item a local gradient reference
    \[
    G^{\mathrm{grad}}_{n,h,j,d}
      =\frac{\partial g_{k,h}(x_n)}
             {\partial x_{j,\tau-d}},
    \]
    computed by automatic differentiation when the sampled operators are
    differentiable and checked against finite differences; and
    \item generator-level Shapley values obtained by evaluating the executable
    generator under feature coalitions and a declared conditional resampling
    rule; when the benchmark declares a hierarchy of player groups, the same
    game also yields Owen values that respect that grouping.
\end{enumerate}

The raw gradient measures local sensitivity. A baseline-relative
gradient--input quantity,
$(x_{j,\tau-d}-b_{j,d})G^{\mathrm{grad}}_{n,h,j,d}$, may also be stored, but it
will be named separately because it is a contribution heuristic rather than a
Shapley value. Comparisons remain within the same explanation family: a model
gradient is compared with the generator gradient, while a model SHAP
explanation is compared with the generator SHAP reference. Exact Shapley
enumeration will validate small cases; permutation sampling and caching will be
used when the number of players makes enumeration impractical.

\subsection{Model-independent evaluation layer}

The public boundary of the benchmark is a \texttt{ForecastAdapter}. Every
adapter declares its input capabilities and implements fitting, prediction,
and, when available, access to an intrinsic mechanism or differentiable model.
The evaluator will therefore distinguish three cases:

\begin{itemize}
    \item every adapter can be evaluated for forecast accuracy;
    \item any callable predictor can receive model-agnostic perturbation or
    finite-difference explanations;
    \item mechanism-recovery metrics are computed only when a method exposes a
    generator-compatible structure, equation, coefficient tensor, or other
    explicit mechanism. Unsupported cells are reported as not applicable, not
    as zero performance.
\end{itemize}

The initial integration set comprises Chronos-2
\citep{ansari2025chronos2}, Toto 2.0 from Datadog
\citep{khwaja2026toto2}, AutoARIMA \citep{hyndman2008auto}, and, as an
optional explainable-by-design baseline, DCIts \citep{baric2026dcits}. The
planning name ``Kronos-2'' is normalized here to \emph{Chronos-2}; Kronos is a
different foundation model specialized for financial-market sequences.
Chronos-2 and Toto 2.0 are evaluated as pretrained forecasters through their
adapters. AutoARIMA is evaluated first with target-only history; a separate
ARIMAX configuration can later receive declared exogenous variables. DCIts
provides a learned transition tensor rather than a symbolic equation, so its
native coefficients can be compared with compatible generator references.

Forecast metrics will include MAE, RMSE, and MASE. Structural recovery will use
precision, recall, and F1 over source--lag support. Coefficient recovery will
use signed error and sign agreement when such coefficients exist. Local
explanation comparisons will include normalized attribution error, rank
correlation, top-$q$ overlap, and sign agreement. All results will retain the
instance, target, horizon, source variable, lag, operator family, formula
complexity, and random seed required for stratified analysis.

\subsection{Validation gates}

The release is accepted only if (i) AST parsing and serialization round-trip
without changing a formula, (ii) seeded generation is reproducible, (iii)
hand-derived support, coefficient, contribution, gradient, and small Shapley
examples match the implementation, (iv) one-step and recursively unrolled
horizon references are separated, (v) all model adapters receive the declared
input only, and (vi) the complete pipeline runs from configuration file to
tables without notebook-only state. These gates protect the benchmark from
producing precise-looking metrics for an incorrectly represented mechanism.

\subsection{Weekly plan}

Table~\ref{tab:weekly-plan} places formula and series generation before the
explanation layer, and places the model study after both. Each week ends in a
testable artifact rather than an open-ended research activity.

Figure~\ref{fig:research-plan-gantt} presents the same plan as a dependency-aware
Gantt chart. Solid arrows mark hard gates. The dashed region identifies work
that can proceed in parallel once the executable generator has passed its first
validation gate; the adapter contract begins earlier from the frozen data
contract, but its implementation converges with the dataset and explanation
branches before model integration.

\begin{landscape}
\begin{figure}[H]
    \centering
    \resizebox{0.98\linewidth}{!}{%
        \begin{tikzpicture}[
    x=1.13cm,
    y=0.70cm,
    tasklabel/.style={font=\scriptsize,anchor=east,align=right,text width=7.0cm},
    bar/.style={draw=black!60,line width=0.35pt,rounded corners=1pt},
    dependency/.style={->,>=stealth,semithick,draw=black!75},
    legend/.style={font=\scriptsize,anchor=west}
]

% Timeline grid.
\fill[black!4] (0,-0.5) rectangle (11,12.5);
\foreach \x in {0,...,11} {
    \draw[black!22,line width=0.25pt] (\x,-0.5) -- (\x,12.5);
}
\foreach \y in {-0.5,0.5,...,12.5} {
    \draw[black!14,line width=0.2pt] (0,\y) -- (11,\y);
}

% Week headings.
\node[font=\scriptsize,align=center] at (0.5,13.05) {21--27\\Sep};
\node[font=\scriptsize,align=center] at (1.5,13.05) {28 Sep--\\4 Oct};
\node[font=\scriptsize,align=center] at (2.5,13.05) {5--11\\Oct};
\node[font=\scriptsize,align=center] at (3.5,13.05) {12--18\\Oct};
\node[font=\scriptsize,align=center] at (4.5,13.05) {19--25\\Oct};
\node[font=\scriptsize,align=center] at (5.5,13.05) {26 Oct--\\1 Nov};
\node[font=\scriptsize,align=center] at (6.5,13.05) {2--8\\Nov};
\node[font=\scriptsize,align=center] at (7.5,13.05) {9--15\\Nov};
\node[font=\scriptsize,align=center] at (8.5,13.05) {16--22\\Nov};
\node[font=\scriptsize,align=center] at (9.5,13.05) {23--29\\Nov};
\node[font=\scriptsize,align=center] at (10.5,13.05) {30\\Nov};

% Task labels.
\node[tasklabel] at (-0.25,12) {Contracts, temporal semantics, and fixed grammar};
\node[tasklabel] at (-0.25,11) {Canonical AST (Week 1) and constrained sampler (Week 2)};
\node[tasklabel] at (-0.25,10) {Series generator, trajectories, and Gaussian innovation};
\node[tasklabel] at (-0.25,9) {Dataset factory and CPU/GPU batched generation};
\node[tasklabel] at (-0.25,8) {Difficulty families and calibration descriptors};
\node[tasklabel] at (-0.25,7) {GT: support, signed coefficients, and local terms};
\node[tasklabel] at (-0.25,6) {GT: gradients and baseline-relative gradient--input};
\node[tasklabel] at (-0.25,5) {GT: Shapley/Owen, exact and permutation-based};
\node[tasklabel] at (-0.25,4) {Universal adapter, capability, and metric interfaces};
\node[tasklabel] at (-0.25,3) {Chronos-2, Toto 2.0, and AutoARIMA integration};
\node[tasklabel] at (-0.25,2) {Full benchmark and failure analysis by mechanism};
\node[tasklabel] at (-0.25,1) {Unit tests, documentation, and reproducibility checks};
\node[tasklabel] at (-0.25,0) {Freeze and archive version 0.1.0};

% Parallel work window after the first validated generator.
\filldraw[fill=orange!4,draw=black!40,dashed,rounded corners=2pt]
    (3.94,3.55) rectangle (8.06,9.45);
\node[font=\scriptsize\itshape,anchor=south west] at (4.05,9.48)
    {parallel work after generator validation};

% Core mechanism construction.
\filldraw[bar,fill=blue!45] (0.05,11.73) rectangle (0.95,12.27);
\filldraw[bar,fill=blue!45] (0.05,10.73) rectangle (1.95,11.27);
\filldraw[bar,fill=blue!45] (2.05,9.73) rectangle (3.95,10.27);

% Scalable generation and controlled difficulty.
\filldraw[bar,fill=green!42] (4.05,8.73) rectangle (7.95,9.27);
\filldraw[bar,fill=green!42] (4.05,7.73) rectangle (7.95,8.27);

% Explanation-ground-truth methods.
\filldraw[bar,fill=orange!58] (4.05,6.73) rectangle (4.95,7.27);
\filldraw[bar,fill=orange!58] (5.05,5.73) rectangle (5.95,6.27);
\filldraw[bar,fill=orange!58] (5.05,4.73) rectangle (6.95,5.27);

% Evaluation interface and model study.
\filldraw[bar,fill=violet!38] (1.05,3.73) rectangle (7.95,4.27);
\filldraw[bar,fill=violet!38] (8.05,2.73) rectangle (8.95,3.27);
\filldraw[bar,fill=violet!38] (9.05,1.73) rectangle (9.95,2.27);

% Continuous quality work and final milestone.
\filldraw[bar,fill=black!28] (0.05,0.73) rectangle (10.95,1.27);
\filldraw[bar,fill=red!48] (10.05,-0.27) rectangle (10.95,0.27);

% Hard dependency gates.
\draw[dependency] (0.95,11.73) -- (0.95,11.45) -- (2.00,11.45)
    -- (2.00,10.20) -- (2.05,10.20);
\draw[dependency] (1.95,11) -- (2.00,11) -- (2.00,9.80) -- (2.05,9.80);
\draw[dependency] (3.95,10) -- (4.00,10) -- (4.00,9.25);
\draw[dependency] (8.00,3.55) -- (8.00,3) -- (8.05,3);
\draw[dependency] (8.95,3) -- (9.00,3) |- (9.05,2);
\draw[dependency] (9.95,2) -- (10.00,2) -- (10.00,0) -- (10.05,0);

% Legend, split over two rows to remain readable after page scaling.
\filldraw[bar,fill=blue!45] (0,-1.45) rectangle (0.35,-1.08);
\node[legend] at (0.47,-1.265) {mechanism};
\filldraw[bar,fill=green!42] (2.60,-1.45) rectangle (2.95,-1.08);
\node[legend] at (3.07,-1.265) {data/difficulty};
\filldraw[bar,fill=orange!58] (5.45,-1.45) rectangle (5.80,-1.08);
\node[legend] at (5.92,-1.265) {explanation GT};
\filldraw[bar,fill=violet!38] (8.40,-1.45) rectangle (8.75,-1.08);
\node[legend] at (8.87,-1.265) {models/evaluation};

\filldraw[bar,fill=black!28] (3.55,-2.00) rectangle (3.90,-1.63);
\node[legend] at (4.02,-1.815) {continuous QA};
\filldraw[bar,fill=red!48] (6.45,-2.00) rectangle (6.80,-1.63);
\node[legend] at (6.92,-1.815) {release};

\end{tikzpicture}
%
    }
    \caption{Dependency-aware Gantt plan for the benchmark release on
    30 November 2026. The explanation-ground-truth branch explicitly covers
    binary source--lag support, signed coefficients and realized local terms,
    automatic-differentiation gradients and gradient--input values, and
    generator-level Shapley values or Owen values for grouped players. Exact
    coalition enumeration is used for small games and permutation estimation
    for larger games.}
    \label{fig:research-plan-gantt}
\end{figure}
\end{landscape}

\begin{landscape}
\begin{figure}[H]
    \centering
    \resizebox{0.77\linewidth}{!}{%
        \begin{tikzpicture}[
    x=1cm,
    y=1cm,
    task/.style={
        draw=black!55,
        line width=0.45pt,
        rounded corners=2pt,
        align=left,
        inner sep=4pt,
        font=\scriptsize
    },
    dependency/.style={->,>=stealth,semithick,draw=black!72},
    joinline/.style={semithick,draw=black!72}
]

% Shared scientific and software foundation.
\node[task,fill=black!6,text width=12.0cm,minimum height=1.25cm] (specification) at (0,7.4) {%
    \textbf{1. Specify semantics, grammar, and public contracts}\par
    Fix $t$, forecast origin $\tau$, horizon $h$, $X_j(t-\ell)$, innovation
    $\varepsilon_t$, the fixed expression vocabulary, configuration fields,
    canonical array axes, and the separation between model inputs and protected metadata.
};

\node[task,fill=blue!8,text width=9.2cm,minimum height=1.45cm] (language) at (-5.2,5.6) {%
    \textbf{2. Implement the canonical mechanism language}\par
    Build immutable AST nodes, the operator registry, validation, canonical JSON,
    formatting, and structural visitors (e.g., represent
    $0.7X_1(t-1)+\sin X_2(t-3)$ and recover its sources and lags).
};

% The adapter boundary can start from the public contracts while generators are built.
\node[task,fill=green!9,text width=9.2cm,minimum height=1.45cm] (adapter) at (5.8,5.6) {%
    \textbf{13. Implement the universal model adapter}\par
    Define capability-aware \texttt{fit}, \texttt{predict}, \texttt{explain}, and
    mechanism-recovery interfaces so any forecaster receives only declared inputs and
    unsupported outputs are reported as not applicable.
};

\node[task,fill=blue!8,text width=7.0cm,minimum height=1.55cm] (sampler) at (-7.4,3.65) {%
    \textbf{3. Sample constrained mechanisms}\par
    Draw unary/binary trees, variables, lags, coefficients, nonlinearities,
    interactions, regimes, and complexity labels from the fixed grammar using
    reproducible seed streams.
};

\node[task,fill=blue!8,text width=7.0cm,minimum height=1.55cm] (compiler) at (0,3.65) {%
    \textbf{4. Compile and screen mechanisms}\par
    Turn each AST into an executable recurrence and reject pilot simulations that
    diverge, become non-finite, collapse to constants, or leave sampled dependencies
    inactive.
};

\node[task,fill=green!9,text width=7.0cm,minimum height=1.55cm] (models) at (7.4,3.65) {%
    \textbf{14. Integrate representative models}\par
    Add adapters for Chronos-2, Toto~2.0, and AutoARIMA, plus optional
    explainable-by-design or symbolic methods such as DCIts when their native
    mechanisms can be exposed.
};

\node[task,fill=blue!8,text width=10.5cm,minimum height=1.45cm] (trajectories) at (0,1.75) {%
    \textbf{5. Generate trajectories from accepted mechanisms}\par
    Sample initial conditions, apply burn-in, recursively generate multivariate and
    multi-target series, create multiple trajectories per formula, and separately
    retain the conditional mean and Gaussian innovation.
};

\node[task,fill=green!9,text width=10.5cm,minimum height=1.45cm] (dataset) at (0,-0.05) {%
    \textbf{6. Build the scalable dataset factory}\par
    Produce chronological windows, splits, horizons, masks, and metadata through lazy
    PyTorch-style batches; benchmark worker-based CPU generation against vectorized
    CPU/GPU generation for large test collections.
};

% Parallel products derived from the retained mechanism and generated observations.
\node[task,fill=violet!8,text width=5.2cm,minimum height=1.65cm] (difficulty) at (-9.0,-2.05) {%
    \textbf{7. Characterize difficulty}\par
    Combine structural descriptors and empirical forecastability (e.g., depth,
    operators, lag span, noise/SNR, stability) to support random and stratified
    easy/medium/hard samples.
};

\node[task,fill=orange!10,text width=5.2cm,minimum height=1.65cm] (structural) at (-3.0,-2.05) {%
    \textbf{8. Derive structural ground truths}\par
    Extract one-step and horizon-unrolled source--lag masks, signed coefficients, and
    realized local terms from the AST and recursive dependency graph.
};

\node[task,fill=orange!10,text width=5.2cm,minimum height=1.65cm] (gradient) at (3.0,-2.05) {%
    \textbf{9. Derive differential ground truths}\par
    Compute per-horizon generator gradients and baseline-relative gradient--input
    values by automatic differentiation and validate smooth cases with finite
    differences.
};

\node[task,fill=orange!10,text width=5.2cm,minimum height=1.65cm] (coalition) at (9.0,-2.05) {%
    \textbf{10. Derive coalition ground truths}\par
    Evaluate generator-level Shapley/Owen games under declared conditional resampling,
    using exact coalitions for small games and cached permutation estimates for larger
    games.
};

\node[task,fill=yellow!12,text width=10.0cm,minimum height=1.55cm] (canonicalapi) at (-5.5,-4.15) {%
    \textbf{11. Assemble the canonical dataset and ground-truth API}\par
    Return history, known covariates, future targets, mechanism/trajectory/window IDs,
    difficulty descriptors, and selected explanation references, always indexed by
    target and forecast horizon.
};

\node[task,fill=yellow!12,text width=10.0cm,minimum height=1.55cm] (evaluation) at (5.5,-4.15) {%
    \textbf{15. Implement capability-aware evaluation}\par
    Score forecasts (MAE/RMSE/MASE), mechanism recovery (support F1, coefficient/sign
    error), and explanations (normalized error, rank, top-$q$, sign); compute only the
    outputs supported by each model/method.
};

% Two intended uses of the generated data and ground truths.
\node[task,fill=green!8,text width=10.0cm,minimum height=1.55cm] (prior) at (-5.5,-6.10) {%
    \textbf{12. Expose a synthetic pretraining-prior interface}\par
    Stream targets for future values, the generating equation/token tree, the
    innovation distribution, and optional explanation references so forecasting,
    symbolic, or explanation-aware models can be pretrained jointly.
};

\node[task,fill=violet!9,text width=10.0cm,minimum height=1.55cm] (benchmark) at (5.5,-6.10) {%
    \textbf{16. Run the benchmark and diagnose failures}\par
    Compare models by horizon, operator, lag, interaction, noise, difficulty, and
    formula complexity to identify which mechanisms are forecast, recovered, or
    explained poorly.
};

\node[task,fill=red!7,text width=12.8cm,minimum height=1.45cm] (release) at (0,-8.10) {%
    \textbf{17. Validate reproducibility, document, and release}\par
    Maintain unit/integration tests throughout development; verify seeds, manifests,
    CPU/GPU agreement, input isolation, and hand-derived examples; publish the library,
    configurations, reproduction commands, reports, and versioned release.
};

% Hard technical dependencies.
\draw[dependency] (specification.south) -- ++(0,-0.22) |- (language.north);
\draw[dependency] (specification.south) -- ++(0,-0.22) |- (adapter.north);

\draw[dependency] (language.south) -- ++(0,-0.22) |- (sampler.north);
\draw[dependency] (language.south) -- ++(0,-0.22) |- (compiler.north);
\draw[dependency] (sampler.east) -- (compiler.west);
\draw[dependency] (compiler.south) -- (trajectories.north);
\draw[dependency] (trajectories.south) -- (dataset.north);

\draw[dependency] (adapter.south) -- ++(0,-0.22) |- (models.north);

\draw[dependency] (dataset.south) -- (difficulty.north);
\draw[dependency] (dataset.south) -- (structural.north);
\draw[dependency] (dataset.south) -- (gradient.north);
\draw[dependency] (dataset.south) -- (coalition.north);

\coordinate (productjoin) at (-0.15,-3.15);
\draw[joinline] (difficulty.south) -- (productjoin);
\draw[joinline] (structural.south) -- (productjoin);
\draw[joinline] (gradient.south) -- (productjoin);
\draw[joinline] (coalition.south) -- (productjoin);
\fill[black!72] (productjoin) circle (1.2pt);
\draw[dependency] (productjoin) -- (canonicalapi.north east);

\draw[dependency] (canonicalapi.south) -- (prior.north);
\draw[dependency] (canonicalapi.east) -- (evaluation.west);
\draw[dependency] (models.east) -- ++(1.2,0) |- (evaluation.north east);
\draw[dependency] (evaluation.south) -- (benchmark.north);

\draw[dependency] (prior.south) -- ++(0,-0.22) |- (release.north west);
\draw[dependency] (benchmark.south) -- ++(0,-0.22) |- (release.north east);

\end{tikzpicture}
%
    }
    \caption{Overall technical dependency graph for the project.  The retained
    symbolic mechanism drives trajectory generation, difficulty characterization,
    and structural, differential, and coalition-based explanation references.  The
    resulting API supports two consumers: synthetic pretraining priors and a
    model-agnostic benchmark with capability-aware forecast, mechanism-recovery,
    and explanation evaluation.}
    \label{fig:project-dependency-graph}
\end{figure}
\end{landscape}

\begingroup
\small
\setlength{\LTpre}{0.4\baselineskip}
\setlength{\LTpost}{0.7\baselineskip}
\renewcommand{\arraystretch}{1.08}
\begin{longtable}{@{}
>{\raggedright\arraybackslash}p{0.14\textwidth}
>{\raggedright\arraybackslash}p{0.25\textwidth}
>{\raggedright\arraybackslash}p{0.28\textwidth}
>{\raggedright\arraybackslash}p{0.25\textwidth}@{}}
\caption{Dependency-aware schedule for a usable benchmark by 30 November 2026.}
\label{tab:weekly-plan}\\
\toprule
Dates & Main work & Deliverable & Acceptance criterion \\
\midrule
\endfirsthead
\multicolumn{4}{@{}l}{\small\emph{Table~\thetable\ continued}}\\
\toprule
Dates & Main work & Deliverable & Acceptance criterion \\
\midrule
\endhead
\midrule
\multicolumn{4}{r@{}}{\small\emph{Continued on the next page}}\\
\endfoot
\bottomrule
\endlastfoot

21--27 Sep. &
Freeze temporal semantics, AST schema, operator registry, and configuration
contract. &
Grammar specification, Python types, and hand-written reference formulas. &
At least 30 representative formulas parse, serialize, and preserve their
source--lag identities. \\

28 Sep.--4 Oct. &
Implement constrained random AST sampling and formula rejection filters. &
Seeded formula sampler with complexity and operator labels. &
A batch of 10,000 candidates is reproducible; every accepted formula is
causal, finite in pilot execution, and reconstructible from metadata. \\

5--11 Oct. &
Compile ASTs and implement initial conditions, burn-in, recursive simulation,
and multivariate output. &
Executable deterministic generator and unit tests. &
Repeated seeds reproduce identical series; simple AR, cross-lag, and nonlinear
examples match hand calculations. \\

12--18 Oct. &
Create dataset splits and controlled families for interactions, irrelevant
variables, lag ranges, and temporal scales. &
Versioned dataset factory and data cards. &
Every sample links back to one formula; leakage checks pass and all information
available to a model is recorded explicitly. \\

19--25 Oct. &
Extract structural support, signed coefficients where defined, realized local
terms, and horizon-unrolled dependencies. &
First ground-truth API and visual sanity reports. &
Hand-derived one-step and multi-step examples match the stored masks and
values exactly. \\

26 Oct.--1 Nov. &
Implement generator gradients and baseline-relative gradient--input outputs. &
Gradient reference module indexed by instance, horizon, variable, and lag. &
Autodiff agrees with central finite differences on smooth formulas; undefined
or non-smooth cases are flagged rather than silently scored. \\

2--8 Nov. &
Implement coalition construction, dependency-aware resampling, exact small
Shapley cases, and permutation estimates. &
Generator SHAP reference module with uncertainty and sampling metadata. &
Exact and sampled results agree within a declared tolerance on small formulas,
and efficiency is checked whenever the chosen game satisfies it. \\

9--15 Nov. &
Implement adapter protocols, forecast/mechanism/explanation metrics, capability
checks, and result schema. &
End-to-end library pipeline using naive and oracle test adapters. &
One configuration produces datasets, references, predictions, explanations,
and stratified tables without notebook state. \\

16--22 Nov. &
Integrate Chronos-2, Toto 2.0, and target-only AutoARIMA; add DCIts if the core
path remains on schedule. &
Model adapters and a common pilot experiment. &
Every model sees the documented context and horizon; unsupported mechanism
metrics are marked not applicable. \\

23--29 Nov. &
Run the comparison matrix; analyze error by operator, lag range, horizon,
interaction, and irrelevant-variable count; finish tests and documentation. &
Release candidate, examples, metric tables, and reproducibility manifest. &
A clean environment can reproduce the principal tables from one command and
all mandatory tests pass. \\

30 Nov. &
Freeze and archive version 0.1.0. &
Installable package, tagged configuration set, generated artifacts, and usage
example. &
The release checklist and a full smoke test pass from installation through
report export. \\
\end{longtable}
\endgroup

\subsection{Final goal: using the benchmark as a library}

The final user should configure a generator once and evaluate multiple models
without changing the data-generation or scoring code. Listing~\ref{lst:final-api}
illustrates the intended public API. The adapter classes belong to the proposed
library; internally, they translate the canonical benchmark batch to the
official input expected by each external package.

\begin{lstlisting}[language=Python,caption={Intended end-state API for generating one benchmark and evaluating the prioritized models.},label={lst:final-api}]
from tsxai_benchmark import Benchmark, GeneratorSpec
from tsxai_benchmark.ground_truth import (
    StructuralSupport,
    SignedCoefficients,
    LocalContributions,
    GeneratorGradient,
    GradientInput,
    GeneratorShapley,
)
from tsxai_benchmark.models import (
    Chronos2Adapter,
    Toto2Adapter,
    AutoARIMAAdapter,
    DCItsAdapter,
)
from tsxai_benchmark.xai import InputGradient, PermutationSHAP

# 1. Sample formulas, execute them, and retain every mechanism.
benchmark = Benchmark.generate(
    GeneratorSpec(
        n_formulas=500,
        series_length=1024,
        context_length=256,
        horizon=24,
        max_lag=48,
        operators=["add", "multiply", "sin", "square"],
        include_irrelevant_variables=True,
        seed=42,
    ),
    ground_truths=[
        StructuralSupport(unroll_horizon=True),
        SignedCoefficients(when_defined=True),
        LocalContributions(),
        GeneratorGradient(),
        GradientInput(baseline="training_mean"),
        GeneratorShapley(
            masker="generator_conditional",
            permutations=512,
        ),
    ],
)

# 2. Model-specific details are isolated inside adapters.
models = [
    Chronos2Adapter(model_id="amazon/chronos-2"),
    Toto2Adapter(model_id="Datadog/Toto-2.0-22m"),
    AutoARIMAAdapter(
        season_length=24,
        input_mode="target_only",
    ),
    # Optional explainable-by-design baseline.
    DCItsAdapter(context_length=48),
]

# 3. Like is compared with like: gradients with gradients,
#    and SHAP values with generator-level Shapley values.
report = benchmark.evaluate(
    models=models,
    explainers={
        "gradient": InputGradient(fallback="finite_difference"),
        "shap": PermutationSHAP(
            masker=benchmark.conditional_sampler,
            permutations=512,
        ),
    },
    forecast_metrics=["mae", "rmse", "mase"],
    mechanism_metrics=["support_f1", "coefficient_error"],
    explanation_metrics=[
        "normalized_error",
        "rank_correlation",
        "top_k_f1",
        "sign_accuracy",
    ],
)

# 4. Diagnose which mechanisms are difficult for each model.
table = report.table(
    groupby=["model", "operator", "lag_band", "horizon"]
)
table.to_csv("results/by_mechanism.csv", index=False)
\end{lstlisting}

The evaluator consults each adapter's capability record. Forecast metrics apply
to all four models. Model-level gradient and SHAP comparisons apply when the
selected explainer can query the fitted predictor. Direct mechanism recovery
applies to DCIts and to future symbolic system-identification adapters that
return an explicit compatible object; it is not imputed for Chronos-2 or Toto
2.0. AutoARIMA's selected orders and coefficients are stored for analysis, but
they are interpreted against the generator only under a declared compatible
linear setting.

A new forecasting model is added by implementing the same narrow contract:

\begin{lstlisting}[language=Python,caption={Minimal extension point for an arbitrary forecasting model.}]
class MyModelAdapter(ForecastAdapter):
    capabilities = {
        "multivariate": True,
        "known_future_covariates": False,
        "native_mechanism": False,
        "autograd": False,
    }

    def fit(self, train_batch):
        self.model.fit(train_batch.history, train_batch.targets)
        return self

    def predict(self, context_batch, horizon):
        return self.model.forecast(context_batch.history, horizon=horizon)

report = benchmark.evaluate(models=[MyModelAdapter()])
\end{lstlisting}

This boundary is the practical final goal: formula generation, dataset
construction, reference explanations, metrics, and stratification remain fixed,
while only the adapter changes. The resulting report can then show not merely
which model forecasts best, but which operators, lag structures, horizons, and
interaction patterns each model uses or fails to recover.

